\subsection{Affine act reductions at the eighteen-arrow frontier}

The following proofs reconstruct the affine reductions at their
points of use. They require no inherited finite-search output.

\begin{lemma}[Proved: a common rational affine row origin]
\label{CP2:A-act-common-affine-origin}
Let $M$ have rectangular minimal ideal $K=I\times\{0,1\}$.
Write $\sigma_m$ and $\tau_m$ for its row and column maps.
Suppose there is $p\in\mathbb QI$, of augmentation one,
fixed by every $\sigma_m$ for which $\tau_m$ is nonconstant.
For every finite nonempty right $M$-act $S$, all products
of rational degree-two classes of $B\mathcal C(M,S)$ vanish.
Neither bounded homology nor a bound on $|S|$ is required.
\end{lemma}
\begin{proof}
First take the actual constant-action quotient of
Lemma~\ref{CP2:A-constant-quotient-H2}. Its degree-two
cohomology surjects onto that of the original cone, so it
suffices to treat constant $K$ action. Compress this cone
by Theorem~\ref{CP2:A-constant-act-compression} to the
monoid $N$ with minimal band $I\times S$. Write $s_j$
for the constant value of an old column $j$; these values
may coincide. They satisfy $s_jm=s_{\tau_m(j)}$.

Choose $i_0\in I$ and use base column $s_0$. Put
$V=\ker(\mathbb QI\to\mathbb Q)$ and
$W=\ker(\mathbb QS\to\mathbb Q)$. The band--Tor
formula identifies $H^2(BN;\mathbb Q)$ with all balanced
forms $L:W\times V\to\mathbb Q$. For any element $f$
of $N$, define
\[
 a_f=[\sigma_f(i_0)]-[i_0],\qquad
 b_f^L(v)=L([\tau_f(s_0)]-[s_0],v),
\]
and let $A_f$ be its row augmentation operator. Balance
and the row/column composition laws give
\[
 a_{fg}=a_f+A_fa_g,\qquad
 b_{fg}^L=b_g^L+b_f^LA_g.
\]
Thus $u_L(f,g)=b_f^La_g$ is the cocycle realizing $L$,
as checked on the band cycles in Lemma~\ref{A:rank-one-square}.

This cocycle transfers through actual functors as follows.
Set $e=(i_0,s_0)$ and use the full Karoubi subcategory
on $1,e$. Extend $u_L$ by the same formula on its
underlying $N$ labels. Both object identities have
$a=b^L=0$, so the extension is normalized. The displayed
identities prove its cocycle equation on every typed pair.
The inclusion of object $1$ is a classifying-space
equivalence, and the compression functor is an equivalence
on classifying spaces. Therefore these pulled-back
cocycles represent all degree-two classes of the act cone.
The compression functor sends a cross state $s$ to
$(i_0,s)$; on that arrow $a_s=A_s=0$ and
$b_s^L=L([s]-[s_0],-)$.

Set $t=p-[i_0]\in V$. Change the row translations on
old monoid arrows to
\[
 a'_m=a_m+(A_m-1)t=\sigma_mp-p.
\]
Set $a'_s=0$ on cross arrows and zero on the sink identity.
These are typed affine coordinates: on any normalized
composable pair, the second arrow is an old monoid arrow
unless the first is the sink identity. Hence the same
affine equation $a'_{fg}=a'_f+A_fa'_g$ holds.
The new cocycle $u'_L(f,g)=b_f^La'_g$ represents the
same class, since
\[
 u'_L-u_L=-\delta h_L,\qquad h_L(f)=b_f^Lt.
\]
For instance the coboundary on a cross/old pair is
$b_f^L(1-A_g)t$, by the displayed identity for $b_{fg}^L$.
All identity cases give zero.

Define, only on the old monoid,
\[
\begin{aligned}
 \chi(m)&=\tau_m(1)-\tau_m(0),&
 \beta(m)&=\tau_m(0),\\
 \ell_L(v)&=L([s_1]-[s_0],v),&
 \alpha_L(m)&=\ell_L(a'_m).
\end{aligned}
\]
Equivariance of the old column line and balance imply
$\ell_LA_m=\chi(m)\ell_L$. Consequently
\[
\begin{aligned}
 \beta(mn)&=\beta(n)+\chi(n)\beta(m),\\
 \alpha_L(mn)&=\alpha_L(m)+\chi(m)\alpha_L(n),\\
 b_m^L&=\beta(m)\ell_L.
\end{aligned}
\]
The common-origin hypothesis is precisely
$\chi(m)a'_m=0$ for every old actor.

Use algebra-ordered normalized chains: a cross arrow,
when present, is first, and every later arrow is old.
For two balanced forms $L,D$, the normalized three-cochain
\[
 P(f,g,h)=u'_L(f,g)\,\beta(g)\,\alpha_D(h)
\]
is therefore defined on every nondegenerate typed triple.
It is zero on degenerate triples. Expanding its five-term
coboundary and substituting the preceding identities gives
\begin{align*}
 (\delta P)(f,g,h,k)
 &=u'_L(f,g)\beta(h)\alpha_D(k)\\
 &\quad+\chi(h)\beta(g)\alpha_D(k)
       \bigl(u'_L(f,gh)-u'_L(f,g)\bigr)\\
 &=u'_L(f,g)u'_D(h,k)
   +\chi(h)\beta(g)\alpha_D(k)b_f^LA_ga'_h\\
 &=u'_L(f,g)u'_D(h,k).
\end{align*}
The last equality uses $\chi(h)a'_h=0$. Thus every
mixed product has a literal rational primitive. Pullback
through the constant-action quotient proves the assertion.
\end{proof}

\begin{lemma}[Proved: projective row dimension at most one]
\label{CP2:A-act-row-pd-one}
Let $M$ have rectangular minimal ideal $I\times J$, put
$A=\mathbb QM/\mathbb QK$ and
$V=\ker(\mathbb QI\to\mathbb Q)$.
If $\operatorname{pd}_A V\leq1$, all rational degree-two
products of every finite right-act cone vanish.
\end{lemma}
\begin{proof}
Take the constant-action quotient and its compression $N$.
The compressed outside algebra is literally $A$, and
its row module is $V$. If $\bar W$ is its column
augmentation module, Lemma~\ref{A:band-Tor} gives
$H_4(BN;\mathbb Q)=\operatorname{Tor}_2^A(\bar W,V)=0$.
The degree-two quotient surjection and naturality of
products give the result, with no boundedness inference.
\end{proof}

\begin{corollary}[Proved: local two-column contracted algebras]
\label{CP2:A-local-two-column-acts}
If $M$ has rectangular minimal ideal with two columns
and its contracted rational algebra is local, all
rational degree-two products of every finite right-act
cone vanish, at any row and state counts.
\end{corollary}
\begin{proof}
A high power of an outside basis element in the pointed
quotient monoid $M/K$ is zero or an idempotent basis
element. A local algebra has only the idempotents zero
and one. In the second case this element is a finite
unit. A nontrivial cyclic group of such units has an
average which is an idempotent different from zero and
one, since its group elements are distinct basis vectors.
Thus every outside element other than the identity is
nilpotent in $A$. The column augmentation character
$\chi:A\to\mathbb Q$ vanishes on them, so their column
maps are constant. Every band actor also has constant
columns. The identity is the only remaining actor, and
it fixes every row. Apply the common-origin lemma.
\end{proof}

\begin{proposition}[Proved: three rows and five outside elements]
\label{CP2:A-three-row-five-outside-acts}
If the rectangular minimal ideal of $M$ has three rows,
two columns, and at most five outside elements, every
finite right-act cone has zero rational degree-two products.
\end{proposition}
\begin{proof}
If $H^2(BM;\mathbb Q)=0$, the suspension quotient of
the monoid vertex already kills all such products.
Otherwise the balanced description of $H^2(BM)$ gives
a nonzero covector $\ell$ on the two-dimensional row
augmentation module satisfying
$\ell A_m=\chi(m)\ell$. Regard its three row values
as a finite rational set, up to a common translation.

If the values are distinct, a slope-one row actor fixes
every row: a nonzero translation cannot preserve a
finite set. Any slope-minus-one actor is its unique
reflection and fixes the middle row. A common row
origin therefore exists. Otherwise translate and rescale
the row values to $(0,0,1)$. A nonconstant column actor
with slope one fixes row2 and preserves the pair of rows0,1.
An actor of slope minus one has one of the two row maps
\[
 r_0=(2,2,0),\qquad r_1=(2,2,1),
\]
and its column map is the swap. If neither occurs, row2
is a common origin. If only $r_0$ occurs, every slope-one
actor fixes row0: composing on the left with $r_0$
would otherwise produce $r_1$. The affine midpoint of
rows0 and2 is then fixed by every nonconstant-column
actor. The other one-flip case is symmetric.
All these cases have zero products by the first lemma.

A possible nonzero product therefore forces both flips,
their squares
$p_0=(0,0,2)$ and $p_1=(1,1,2)$ with identity columns,
and the actual identity. These five distinct profiles
are all outside the minimal ideal. The outside budget
is therefore exactly five, with one element of each
profile. Products among them stay outside, since column
maps remain bijective. Their row and column profiles
force the multiplication
\[
 (i,\epsilon)(j,\delta)=(i,\epsilon+\delta),
 \qquad (i,0)=p_i,\quad(i,1)=r_i,
 \quad\epsilon,\delta\in\mathbb Z/2.
\]
The contracted algebra is consequently the unitalization
of $\mathbb Q[L_2\times C_2]$.

Put $e_-=(p_0-r_0)/2$ and
$n_-=(p_1-p_0-r_1+r_0)/2$. Direct multiplication gives
$e_-^2=e_-$ and $Ae_-=\mathbb Qe_-\oplus\mathbb Qn_-$.
In the row basis $u=[1]-[0]$, $v=[2]-[0]$ we have
$e_-v=v$ and $n_-v=-u$. The $A$-linear map
$Ae_-\to V$, $ae_-\mapsto av$, is therefore an
isomorphism. Thus $V$ is projective and the preceding
row-dimension lemma excludes the last possibility.
This proves the assertion without a source catalog.
\end{proof}

\begin{proposition}[Proved: two columns and three outside elements]
\label{CP2:A-three-outside-two-column-acts}
If the rectangular minimal ideal has two columns and at
most three outside elements, every finite right-act cone
has zero rational degree-two products, at any row count.
\end{proposition}
\begin{proof}
Let $P$ be the outside actors with nonconstant column
map. Their slopes multiply in $\{1,-1\}$, so $P$ is
an actual monoid containing the identity, of order at
most three. If $P$ is a group, average a row orbit.
If its minimal ideal is a singleton zero $z$, any row
in the image of $\sigma_z$ is fixed by $P$. The only
other possible proper minimal ideal has two elements:
it is a group of order two, a left-zero band, or a
right-zero band. A group ideal has a $P$-invariant
affine average, since the only outside element is the
identity. For a right-zero ideal, $xe=e$ shows that
every row in the image of $\sigma_e$ is fixed by $P$.
The common-origin lemma applies in all these cases.

The remaining monoid is the identity-adjoined left-zero
two-band. It uses all three outside elements, so
$A=\mathbb QP$. If its two left zeros are $p,q$, the
orthogonal idempotents $p,1-p$ and the vector $q-p$
give the three-dimensional one-arrow path algebra:
the radical squares to zero and lies in $(1-p)Ap$.
The one-arrow path algebra is hereditary. For example
its representation consists of a single linear map;
splitting that map into kernel, image and cokernel
decomposes every module into the two simples and the
single arrow module, whose projective resolutions have
length at most one. Hence $\operatorname{pd}_A V\leq1$,
and the row-dimension lemma applies.
\end{proof}

\begin{theorem}[Proved: bounded act cones through order eleven]
\label{CP2:A-bounded-eleven-acts}
For any finite monoid $M$ of order at most eleven and
any finite nonempty right act $S$, bounded rational
homology of $B\mathcal C(M,S)$ forces all rational
degree-two products to vanish. There is no bound on
the state count or the second Betti number. The left-act
version and the exact-projective-opposite-corner version
follow by opposition and actual orbit comparison.
\end{theorem}
\begin{proof}
If the monoid vertex has zero rational second cohomology,
use its suspension quotient. At most two minimal rows
are excluded by Lemma~\ref{CP2:A-high-two-row-act}.
Otherwise there are at least three rows and two columns.
A nontrivial maximal group would require at least
$3\cdot2\cdot2+1=13$ elements, so the ideal is a band.
Writing $q=|M\setminus K|$, the complete size list is
\[
\begin{gathered}
 (r,c,q)=(3,2,q\leq5),\quad (4,2,q\leq3),\\
 (5,2,q\leq1),\quad (3,3,q\leq2).
\end{gathered}
\]
The first case is the three-row proposition. The next
two are the three-outside proposition. In the last case
Lemma~\ref{CP2:A-two-outside-acts}, applied to the reduced
augmentation module of the actual act, gives zero fourth
rational homology from boundedness. This proves the theorem.
For an opposite projective endpoint, its exact coinvariant
functor is the actual finite orbit quotient and preserves
the rational homology and cohomology ring, including
boundedness. The same conclusion therefore transfers.
\end{proof}

\begin{corollary}[Proved: projective initial singleton through order eleven]
\label{CP2:A-eleven-projective-singleton}
In a full-linear three-object category, suppose the initial
corner has projective trivial left rational module, the
initial adjacent right act has one orbit, and the final
corner has order at most eleven. Bounded rational homology
forces all rational degree-two products to vanish.
The opposite assertion holds at the final endpoint.
\end{corollary}
\begin{proof}
The actual exact-orbit quotient makes the initial corner
trivial and its adjacent hom-set singleton. Its singleton
comma category deletes the middle object by an actual
homotopy equivalence. The remaining ordinary act cone
over the final corner has bounded rational homology.
Apply the preceding theorem and pull back its product
vanishing along the specified comparisons.
\end{proof}
